\documentclass{article}

% if you need to pass options to natbib, use, e.g.:
%     \PassOptionsToPackage{numbers, compress}{natbib}
% before loading neurips_2026

% The authors should use one of these tracks.
% Before accepting by the NeurIPS conference, select one of the options below.
% 0. "default" for submission
\usepackage[education, final]{neurips_2026}
% the "default" option is equal to the "main" option, which is used for the Main Track with double-blind reviewing.
% 1. "main" option is used for the Main Track
%  \usepackage[main]{neurips_2026}
% 2. "position" option is used for the Position Paper Track
%  \usepackage[position]{neurips_2026}
% 3. "eandd" option is used for the Evaluations & Datasets Track
 % \usepackage[eandd]{neurips_2026}
% 4. "creativeai" option is used for the Creative AI Track
%  \usepackage[creativeai]{neurips_2026}
% 5. "sglblindworkshop" option is used for the Workshop with single-blind reviewing
 % \usepackage[sglblindworkshop]{neurips_2026}
% 6. "dblblindworkshop" option is used for the Workshop with double-blind reviewing
%  \usepackage[dblblindworkshop]{neurips_2026}
% 7. "education" option is used for the Education Track
%  \usepackage[education]{neurips_2026}


% After being accepted, the authors should add "final" behind the track to compile a camera-ready version.
% 1. Main Track
 % \usepackage[main, final]{neurips_2026}
% 2. Position Paper Track
%  \usepackage[position, final]{neurips_2026}
% 3. Evaluations & Datasets Track
 % \usepackage[eandd, final]{neurips_2026}
% 4. Creative AI Track
%  \usepackage[creativeai, final]{neurips_2026}
% 5. Workshop with single-blind reviewing
%  \usepackage[sglblindworkshop, final]{neurips_2026}
% 6. Workshop with double-blind reviewing
%  \usepackage[dblblindworkshop, final]{neurips_2026}
% Note. For the workshop paper template, both \title{} and \workshoptitle{} are required, with the former indicating the paper title shown in the title and the latter indicating the workshop title displayed in the footnote.
% For workshops (5., 6.), the authors should add the name of the workshop, "\workshoptitle" command is used to set the workshop title.
% \workshoptitle{WORKSHOP TITLE}
% 7. Education Track
% \usepackage[education, final]{neurips_2026}

% "preprint" option is used for arXiv or other preprint submissions
 % \usepackage[preprint]{neurips_2026}

% to avoid loading the natbib package, add option nonatbib:
%    \usepackage[nonatbib]{neurips_2026}

\usepackage[utf8]{inputenc} % allow utf-8 input
\usepackage[T1]{fontenc}    % use 8-bit T1 fonts
\PassOptionsToPackage{hyphens}{url} % let long URLs break across lines
\usepackage{hyperref}       % hyperlinks
\usepackage{url}            % simple URL typesetting
\usepackage{booktabs}       % professional-quality tables
\usepackage{amsfonts}       % blackboard math symbols
\usepackage{nicefrac}       % compact symbols for 1/2, etc.
\usepackage{microtype}      % microtypography
\usepackage{xcolor}         % colors

% Note. For the workshop paper template, both \title{} and \workshoptitle{} are required, with the former indicating the paper title shown in the title and the latter indicating the workshop title displayed in the footnote. 
\title{Speculative Decoding Tutorial}


% The \author macro works with any number of authors. There are two commands
% used to separate the names and addresses of multiple authors: \And and \AND.
%
% Using \And between authors leaves it to LaTeX to determine where to break the
% lines. Using \AND forces a line break at that point. So, if LaTeX puts 3 of 4
% authors names on the first line, and the last on the second line, try using
% \AND instead of \And before the third author name.


\author{%
  Rachel L. Thomas\\
  AnswerAI \\
  \texttt{rt@answer.ai} 
  %
}


\begin{document}


\maketitle


\begin{abstract}
Speculative decoding is a useful way of accelerating inference in autoregressive models \citep{leviathan_fast_2023, chen_accelerating_2023}. Introduced less than 4 years ago, it has sparked an active line of research, including the algorithms Lookahead, Medusa, and EAGLE \citep{fu_break_2024, cai_medusa_2024, li_eagle_2024}. Speculative decoding is now deployed in many production systems, such as TensorRT-LLM and Google Search AI summary \citep{shah_boost_2024, leviathan_looking_2024}. Recent work continues, such as DeepSeek's DSpark \citep{cheng_dspark_2026}. This tutorial gives an understanding of what the algorithm does through an interactive visualization, provides an empirical check on correctness using a toy data set, and uses a simulation to check the speedup the paper predicts.

\end{abstract}

\section{Leveling \& prerequisites}

This tutorial is aimed at advanced undergraduates. It assumes familiarity with python, probability distributions, and a conceptual understanding of how LLMs generate text using tokens. It does not assume any prior exposure to speculative decoding, to inference optimization, or to the internals of transformer architectures.

\section{Learning objectives \& outcomes}

The tutorial has three objectives. The first is understanding how the algorithm works, using an interactive visualization. The second is to check correctness empirically on a toy data set, to see that speculative decoding results in the same probability distribution as the original, slower model. The third is to see where the speedup comes from, by checking the paper's formula for expected tokens per iteration against a Monte Carlo simulation.

Readers who works through the tutorial will be able to explain the key assumption of memory bandwidth being much more expensive than arithmetic computation, and why this makes a batch of guessed tokens almost as cheap to verify as a single one. They will also be able to explain why a rejected token is resampled from the residual distribution rather than from the target distribution. They will be able to describe the three quantities the speedup depends on, namely the acceptance rate, the number of draft tokens per iteration, and the cost of the draft model relative to the target, and explain why improving the draft model pulls two of them in opposite directions.

\section{Linked papers}

Speculative decoding was introduced independently by \citet{leviathan_fast_2023} and \citet{chen_accelerating_2023}, and it has since become an active line of research. Medusa \citep{cai_medusa_2024}, Lookahead decoding \citep{fu_break_2024}, and EAGLE \citep{li_eagle_2024} were 3 new variations of speculative decoding introduced at ICML 2024. Each revisits a different part of the original design: how the guesses are produced, whether a separate draft model is needed at all, and what the drafter should predict. Research continues, with DSpark \citet{cheng_dspark_2026}, a speculative decoding variation that drafts a block of tokens at once and adaptively adjusts how many guesses get checked.

The method is also used in large production systems. \citet{leviathan_looking_2024} describe its deployment in Google Search AI summaries, and \citet{shah_boost_2024} reports how NVIDIA TensorRT-LLM saw a 3-fold throughput improvement on a large language model when using speculative decoding.

\section{Teaching Materials}

The teaching materials are a Jupyter notebook, \texttt{speculative-decoding-tutorial.ipynb}, together with the supporting source code it imports.

The notebook opens with an interactive visualization of speculative decoding, built with FastHTML and HTMX and served from inside the notebook itself. Each step shows the small model's guesses, which of those guesses the large model accepts, and where the first rejection falls. A screen recording is included for readers who are unable to run the code.

The notebook then works through Algorithm 1 on a toy vocabulary of six tokens. It draws 100,000 samples and compares the empirical distribution against the target distribution. It then breaks the algorithm in two plausible ways, first by resampling from the target distribution instead of from the residual, and second by accepting every proposed token. Both broken versions visibly diverge from the target distribution, while the correct one does not.

The final part turns to speed. It presents the paper's formula for the expected number of tokens per iteration and checks it against a Monte Carlo simulation of the acceptance process. It then introduces the cost of the draft model relative to the target, and the walltime speedup formula that follows, along with the tension this creates: a draft model that gets accepted more often also costs more to run. The notebook closes with an exercise asking the reader to extend the simulation to walltime and compare the result against the published measurements.

A second notebook, \texttt{specdec-viz-implementation.ipynb}, documents how the visualization was built, for readers who want to modify it. Its source is included as \texttt{specdec\_viz\_implementation.py}.

All code runs on a laptop. No GPU, model download, nor API key is required.

\begin{ack}
This work received no specific funding. The author declares no competing interests.
\end{ack}

\medskip


{
\small

\bibliographystyle{plainnat}
\bibliography{ref}

}





\end{document}
